\documentclass[11pt]{article}
% ===========================================================================
%  Shared preamble for the AP Calculus AB exam-preparation worksheets.
%  Input right after \documentclass. Files beginning with "_" are not compiled
%  on their own. Navy and gold, to match the course's lesson pages.
% ===========================================================================
\usepackage[letterpaper,top=2.3cm,bottom=2.2cm,left=1.7cm,right=1.7cm,headheight=28pt,footskip=30pt]{geometry}
\usepackage{amsmath,amssymb}
\usepackage[T1]{fontenc}
\usepackage[utf8]{inputenc}
\usepackage{lmodern}
\usepackage{xcolor}
\usepackage[most]{tcolorbox}
\usepackage{enumitem}
\usepackage{fancyhdr}
\usepackage{array}
\usepackage{tikz}
\usepackage{pgfplots}
\pgfplotsset{compat=1.18}

\definecolor{navy}{HTML}{161E54}
\definecolor{ink}{HTML}{27357E}
\definecolor{gold}{HTML}{B98A00}
\definecolor{goldbg}{HTML}{FFFDF3}
\definecolor{inkbg}{HTML}{F6F7FD}
\definecolor{rule}{HTML}{C9CFF0}
\definecolor{okgreen}{HTML}{1B7A44}
\definecolor{okbg}{HTML}{F1F8F3}

\pagestyle{fancy}
\fancyhf{}
\fancyhead[L]{\footnotesize Name: \underline{\hspace{4cm}}}
\fancyhead[C]{\textbf{\color{navy}Integration Academy}\\[-2pt]{\scriptsize Calculus AB \textperiodcentered{} AP\textsuperscript{\textregistered} Exam Preparation}}
\fancyfoot[L]{\tiny\color{black!55}\parbox[t]{0.78\textwidth}{AP\textsuperscript{\textregistered} is a registered trademark of the College Board, which was not involved in the production of, and does not endorse, this material. \textcopyright{} Integration Academy.}}
\fancyfoot[R]{\scriptsize\thepage}
\renewcommand{\headrulewidth}{1.2pt}
\renewcommand{\headrule}{\hbox to\headwidth{\color{gold}\leaders\hrule height \headrulewidth\hfill}}
\renewcommand{\footrulewidth}{0pt}
\setlength{\parindent}{0pt}
\setlength{\parskip}{4pt}

% Title block: \wstitle{code}{title}{unit line}
\newcommand{\wstitle}[3]{\begin{center}{\small\bfseries\color{gold}WORKSHEET #1}\\[3pt]{\LARGE\bfseries\color{navy}#2}\\[3pt]{\small\color{black!60}#3}\end{center}\vspace{2pt}}
% Section heading
\newcommand{\sech}[1]{\par\vspace{10pt}{\large\bfseries\color{navy}#1}\par\nobreak\vspace{2pt}{\color{rule}\hrule height0.8pt}\vspace{7pt}}
% A level tag printed beside a group of questions
\newcommand{\level}[1]{\par\vspace{6pt}{\footnotesize\bfseries\color{gold}\MakeUppercase{#1}}\par\nobreak\vspace{2pt}}

\newtcolorbox{ideabox}{colframe=ink,colback=inkbg,coltitle=white,colbacktitle=ink,fonttitle=\bfseries\small,title={Key ideas},arc=2pt,boxrule=0.8pt,left=8pt,right=8pt}
\newtcolorbox{example}[1]{colframe=rule,colback=white,coltitle=navy,colbacktitle=inkbg,fonttitle=\bfseries\small,title={#1},arc=2pt,boxrule=0.8pt,borderline west={3pt}{0pt}{ink},left=10pt,right=8pt}
\newtcolorbox{frbox}[1]{colframe=gold,colback=goldbg,coltitle=white,colbacktitle=gold,fonttitle=\bfseries\small,title={#1},arc=2pt,boxrule=0.8pt,breakable,left=8pt,right=8pt}
\newtcolorbox{solbox}{colframe=okgreen,colback=okbg,coltitle=white,colbacktitle=okgreen,fonttitle=\bfseries\small,title={Worked solutions},arc=2pt,boxrule=0.8pt,breakable,left=8pt,right=8pt}
\newtcolorbox{scorebox}[1]{colframe=navy,colback=white,coltitle=white,colbacktitle=navy,fonttitle=\bfseries\small,title={#1},arc=2pt,boxrule=0.8pt,breakable,left=8pt,right=8pt}

\newcommand{\soln}{\par\smallskip\textbf{\color{ink}Solution.}\ }
\newcommand{\result}[1]{\par\smallskip\textbf{#1}}
% blank working space of a given height
\newcommand{\work}[1]{\par\vspace{#1}}
% a point value in the margin of a scoring guide
\newcommand{\pts}[1]{\hfill{\footnotesize\bfseries\color{gold}[#1]}}

% Coordinate grid for graph questions: \begin{apgraph}{xmin}{xmax}{ymin}{ymax} ... \end{apgraph}
\newenvironment{apgraph}[4]{\begin{center}\begin{tikzpicture}\begin{axis}[
  axis lines=middle, xmin=#1, xmax=#2, ymin=#3, ymax=#4,
  grid=both, grid style={black!12}, xtick distance=1, ytick distance=1,
  ticklabel style={font=\scriptsize}, xlabel={$x$}, ylabel={$y$},
  width=10cm, height=7.2cm, samples=80, clip=false,
  every axis plot/.append style={line width=1.1pt, color=ink}]}{\end{axis}\end{tikzpicture}\end{center}}
\newcommand{\closeddot}[2]{\addplot[only marks, mark=*, mark size=2.2pt, color=ink] coordinates {(#1,#2)};}
\newcommand{\opendot}[2]{\addplot[only marks, mark=*, mark size=2.2pt, color=ink, mark options={fill=white}] coordinates {(#1,#2)};}

% Multiple-choice options: four across, or stacked when the options are long.
\newcommand{\opts}[4]{\par\vspace{3pt}\begin{tabular}{@{}p{0.235\linewidth}p{0.235\linewidth}p{0.235\linewidth}p{0.235\linewidth}@{}}(A)\ #1 & (B)\ #2 & (C)\ #3 & (D)\ #4\end{tabular}\par\vspace{9pt}}
\newcommand{\optsv}[4]{\par\vspace{3pt}\begin{tabular}{@{}l@{\ \ }p{0.9\linewidth}@{}}(A) & #1\\(B) & #2\\(C) & #3\\(D) & #4\end{tabular}\par\vspace{9pt}}

\fancyhead[R]{\footnotesize\color{black!60}Worksheet 9.1}
\begin{document}

\wstitle{9.1}{How the Exam Is Built and Scored}{Unit 9: Exam Strategy and Free-Response Types}

\begin{ideabox}
\begin{itemize}[leftmargin=*,itemsep=2pt,topsep=2pt]
  \item \textbf{Section I}: 29 multiple-choice questions in 62 minutes without a calculator, then 13 in 38 minutes with one.
  \item \textbf{Section II}: 2 free-response questions in 30 minutes with a calculator, then 4 in 60 minutes without one.
  \item Each multiple-choice question is worth 1 point and each free-response question 9. The two sections count equally. A wrong multiple-choice answer costs nothing.
  \item On free response, the set-up earns points. Write the integral, equation or derivative before its value.
\end{itemize}
\end{ideabox}

\sech{Worked examples}
\begin{example}{Example 1 \textperiodcentered{} Pacing}
In the calculator multiple-choice part you have answered 9 of the 13 questions and 10 minutes remain. Are you on pace?\soln 4 questions in 10 minutes is 2.5 minutes each; the average for the part is about 2.9.\result{Slightly behind. Answer what you can quickly and leave no question blank.}
\end{example}

\begin{example}{Example 2 \textperiodcentered{} Combining the sections}
A student earns 34 of 42 multiple-choice points and 27 of 54 free-response points. What fraction of the total is this?\soln $\dfrac12\left(\dfrac{34}{42}+\dfrac{27}{54}\right)$.\result{About 0.655, or 65.5\%.}
\end{example}

\begin{example}{Example 3 \textperiodcentered{} Decimal answers}
The value of an integral is $3.14159\ldots$ Which of $3.1$, $3.141$ and $3.142$ would be accepted?\soln A decimal answer must be correct to three places after the point, rounded or truncated.\result{$3.141$ and $3.142$, but not $3.1$.}
\end{example}

\newpage
\sech{Practice}

\level{Core}
\begin{enumerate}[leftmargin=*,itemsep=2pt]
  \item How many free-response questions are there, and on how many of them is a calculator used?\work{2.3cm}
  \item How long is Section I in total?\work{2.3cm}
  \item What is the greatest number of points available in Section II?\work{2.3cm}
  \item What is the average time per question in the no-calculator multiple-choice part?\work{2.3cm}
\end{enumerate}

\level{Exam level}
\begin{enumerate}[leftmargin=*,itemsep=2pt,start=5]
  \item A student earns 25 of 42 multiple-choice points and 40 of 54 free-response points. What fraction of the total is this?\work{2.8cm}
  \item A student expects 36 of 54 on the free response and wants at least 60\% overall. How many multiple-choice questions must be correct?\work{2.8cm}
  \item With 20 minutes left in the no-calculator free-response part, two questions are untouched. Describe a sensible plan.\work{2.8cm}
  \item In part (b) you need your answer to part (a), which you suspect is wrong and cannot fix. What should you do?\work{2.8cm}
\end{enumerate}

\begin{frbox}{Question 9 \textperiodcentered{} Free-response practice}
\emph{No calculator. Write each answer as you would in the exam: set-up first, then the value, then the reason.} Let $f(x)=x^3-3x^2$.
\begin{enumerate}[label=(\alph*),leftmargin=*,itemsep=2pt]
  \item Find the average value of $f$ on $[0,2]$.\work{2.6cm}
  \item Write an equation of the line tangent to the graph of $f$ at $x=1$.\work{2.6cm}
  \item Show that $f$ has a relative minimum at $x=2$. Justify your answer.\work{2.6cm}
\end{enumerate}
\end{frbox}

\newpage
\begin{solbox}
\begin{enumerate}[leftmargin=*,itemsep=4pt]
  \item 6 questions; a calculator is used on 2.
  \item $62+38=100$ minutes.
  \item $6\times9=54$.
  \item $\dfrac{62}{29}\approx2.1$ minutes.
  \item $\dfrac12\left(\dfrac{25}{42}+\dfrac{40}{54}\right)\approx0.668$.
  \item $\dfrac12\left(\dfrac{x}{42}+\dfrac{36}{54}\right)\ge0.6$ gives $x\ge22.4$: at least 23.
  \item About 10 minutes each. Read every part of both, write the set-ups you can, and attempt the parts that stand alone first.
  \item Use your own value and carry on. Work that follows correctly from an earlier error can still earn full credit.
  \item \begin{enumerate}[label=(\alph*),leftmargin=*,itemsep=3pt]
    \item $\dfrac12\displaystyle\int_{0}^{2}(x^3-3x^2)\,dx=\dfrac12\Big[\tfrac{x^4}{4}-x^3\Big]_0^2=\dfrac12(-4)=-2$.
    \item $f(1)=-2$ and $f'(1)=3-6=-3$: $y+2=-3(x-1)$.
    \item $f'(x)=3x(x-2)$, so $f'(2)=0$, and $f'$ changes from negative to positive at $x=2$. Therefore $f$ has a relative minimum there.
  \end{enumerate}
\end{enumerate}
\end{solbox}

\end{document}
